% ARCHIVED AUDITED DRAFT: incorporated into part_B.tex 2026-11-19;
% independent root targeted proof audit PASS 14:36 UTC.
\begin{saresult}{B.044}{Proved}{Integer-valued polynomials commute with one-dimensional finite-residue completion}
Let $(R,\mathfrak m)$ be a one-dimensional Noetherian local
domain with finite residue field, let $K=\operatorname{Frac}(R)$,
and put $A=\widehat R$ and $Q=K\otimes_R A$.
For every number of variables $n$ there is a canonical isomorphism
\[
 \operatorname{Int}(R^n,R)\otimes_R A
       \simeq \operatorname{Int}(A^n,A)\subset Q[X_1,\ldots,X_n].
\]
It identifies the tensor maps after faithfully flat base change.
In particular a noninjective second tensor map for $A\subset Q$
implies a noninjective second tensor map for $R\subset K$.
\end{saresult}
\begin{proof}
Fix bounds on the degree in each variable and let $N$ be the
number of allowed monomials. Choose a rectangular grid of
distinct points in $R$. The product of its Vandermonde
determinants gives a nonzero $d\in R$ such that every
integer-valued polynomial within these bounds has all
coefficients in $d^{-1}R$. The same grid over $A$ shows that
all coefficients there lie in $d^{-1}A$: flatness makes $d$
a nonzerodivisor of $A$.

The ring $R/dR$ is a finite ring, because it is Artinian local
with finite residue field. Choose a finite set $E\subset R$
of representatives. For $f\in d^{-1}R[X_1,\ldots,X_n]$,
its value modulo $R$ depends only on each input modulo $dR$.
Indeed $df$ has coefficients in $R$, and changing a coordinate
by a multiple of $d$ changes its numerator by a multiple of
$d$. Consequently the integer-valued module in these bounds
is the kernel of the finite evaluation map
\[
 (d^{-1}R)^N\longrightarrow (d^{-1}R/R)^{E^n}.
\]
The natural map $R/dR\to A/dA$ is an isomorphism: the source
has finite length and hence is already complete. Thus exactly
the same finite tests describe the bounded-degree module over
$A$. Flat base change commutes with this kernel and identifies
the two modules. Taking the filtered union of all degree
bounds proves the assertion. Polynomial multiplication in
disjoint variables commutes with these identifications, so
the tensor maps correspond as stated. Faithful flatness
detects nonzero kernels.

For clarity, $Q$ is the total quotient ring of $A$. Choose
$0\ne t\in\mathfrak m$. Every $0\ne r\in R$ has
$t^M\in rR$ for some $M$, so $Q=A[1/t]$. This is a
zero-dimensional Noetherian ring: $A/tA$ is Artinian, and
completion has dimension one. It is Artinian and is a
finite product of local Artinian rings. Every nonzerodivisor
of $A$ becomes a unit in this product, so $A[1/t]$ is indeed
the total quotient ring.
The completion facts invoked here are the flatness and exactness
of finite-module completion (Stacks, Lemma 10.97.2, tag 00MB)
and its local faithful-flat consequence.
\end{proof}

\begin{saresult}{B.045}{Proved}{A generic nilpotent forces a tensor kernel after completion}
Let $(A,\mathfrak m)$ be a one-dimensional complete
Noetherian local ring with finite residue field. Assume that
$A$ is Cohen--Macaulay and that either $A$ has characteristic
$p>0$, or $A$ has characteristic zero and $p$ is a
nonzerodivisor, where $p$ is its residue characteristic.
If $A$ is not reduced, then
\[
 \operatorname{Int}(A)\otimes_A\operatorname{Int}(A)
       \longrightarrow\operatorname{Int}(A^2)
\]
is not injective. Here coefficients belong to the total
quotient ring of $A$. The proof constructs a nonzero tensor
and an explicit balanced linear functional detecting it.
\end{saresult}
\begin{proof}
We first construct the finite lattice in which the argument
will take place, including the needed compatible coefficient
field. There is a complete discrete valuation subring
$V_0\subset A$ such that $A$ is finite and torsion-free over
$V_0$, and such that some nonreduced local factor $Q_0$ of
$Q(A)$ has residue field $F$ separable over
$K_0=\operatorname{Frac}(V_0)$.

In mixed characteristic take the Cohen coefficient ring.
It is a complete discrete valuation ring with uniformizer
$p$, and $p$ is a parameter. The finite-module assertion
follows by lifting finitely many generators modulo $p$ and
successively approximating in the $p$-adic topology;
$A$ is $p$-adically complete because $pA$ is
$\mathfrak m$-primary. Separability holds in characteristic
zero.

Here is the extra parameter choice needed in characteristic
$p$. The roots of $X^{|A/\mathfrak m|}-X$ give a coefficient
field $k\subset A$ (or use the Cohen structure theorem).
Choose a minimal prime $\mathfrak p$ corresponding to a
nonreduced generic factor. Such a factor exists because a
parameter is a nonzerodivisor, so a nonzero nilpotent remains
nonzero in $Q(A)$. The complete one-dimensional domain
$C=A/\mathfrak p$ has finite normalization, a complete
discrete valuation ring $V=l[[u]]$ with $l/k$ finite.
Here finiteness follows from the Nagata property of complete
Noetherian local rings; a finite domain over a complete local
ring is local, because a finite complete algebra decomposes
as a product of complete local algebras. A nonzero conductor
then gives $u^N\in C$ for all sufficiently large $N$.
Choose such an $N$ with $p\nmid N$ and lift $u^N$ to
$t_0\in\mathfrak m$.

We can adjust this lift to a parameter without changing its
image in $C$. Choose
$a\in\mathfrak p$ outside every other minimal prime, by
prime avoidance. At each other minimal prime, the nonzero
images of $a,a^2,a^3,\ldots$ are distinct: $a$ lies in
the maximal ideal of the corresponding local domain.
Thus $t_0+a^j$ belongs to each such prime for at most one
$j$. Choosing $j$ outside this finite exceptional set gives
a parameter $t$ mapping to $u^N$. If there is no other
minimal prime, use $t=t_0$.
Set $V_0=k[[t]]\subset A$. This map is injective, since
a nonzero power series is $t^r$ times a unit and $t$ is a
nonzerodivisor. Successive approximation modulo $t^m$ shows
that $A$ is finite over $V_0$. In the chosen factor the
residue field is $F=l((u))$, a finite separable extension
of $k((u^N))$ since $p\nmid N$.

Since $F/K_0$ is finite separable, it embeds in $Q_0$ as
a coefficient field compatibly with $K_0$. To see this
directly, choose a primitive element of $F/K_0$ and lift
it to $Q_0$. Its separable minimal polynomial has invertible
derivative. Newton correction modulo successive powers of
the nilradical, which is nilpotent, gives an actual root
lifting the primitive element. Its $K_0$-algebra is the
required copy of $F$. Write
\[
 \rho:Q_0\longrightarrow F,\qquad N=\ker\rho,
 \qquad N^h=0\quad(h\ge2).
\]
Let $V\subset F$ be the integral closure of $V_0$ in $F$.
It is a finite complete discrete valuation ring, with
uniformizer $u$ and finite residue field. The same notation
$V=l[[u]]$ applies in the equal-characteristic construction.
The ring
\[
 O=V[\operatorname{pr}_0(A)]\subset Q_0
\]
is finite and torsion-free over $V$, contains $V$, satisfies
$\rho(O)=V$, and spans $Q_0$ over $F$.
Finiteness follows because $\operatorname{pr}_0(A)$ is
finite over $V_0$ and hence integral over $V$.
Let $e$ denote the idempotent of the selected factor in
$Q(A)$. Since $eO$ is a finite $V_0$-module in $Q(A)$,
clearing finitely many denominators gives an integer $c\ge0$
such that
\[
                         u^c eO\subset A.                 \tag{B45.1}
\]
Indeed a sufficiently high power of the uniformizer of
$V_0$ works, and in $V$ that power is a unit times a
power of $u$; the unit preserves $O$.

Put $J=N\cap O$. There is a primitive element
$z_0\in J\cap\operatorname{Ann}_{Q_0}(N)$: start with a
nonzero element in the last nonzero power of $N$, multiply
it into $O$, and divide by the greatest possible power of
$u$ in the finite free $V$-module $O$. Choose a
$V$-linear map $\ell:O\to V$ with $\ell(z_0)=1$ and
extend it $F$-linearly to $Q_0$. Such a map exists because
a primitive vector in a finite free module over a discrete
valuation ring extends to a basis.

There is an integer $b\ge0$, independent of
$G\in\operatorname{Int}(A)$, such that, with
$x_G=\operatorname{pr}_0(G'(0))$,
\[
                  Jx_G\subset u^{-b}O,\qquad
                  \rho(x_G)\in u^{-b}V.                   \tag{B45.2}
\]
Here is a finite interpolation proof. Choose finitely many
$V$-generators $n_1,\ldots,n_s$ of $J$. By (B45.1),
$u^c e n_j\in A$. For any $h$ fixed distinct scalars
$s_1,\ldots,s_h\in V_0$, the values
$G(s_i u^c e n_j)$ lie in $A$ and their selected components
lie in $O$. As a polynomial in the scalar $s$, this
evaluation has degree at most $h-1$, because $n_j^h=0$.
Inverting this one fixed Vandermonde matrix bounds its
linear coefficient $u^c n_jx_G$ in a fixed fractional
lattice $u^{-d}O$. Increasing $b$ gives the first inclusion
for every generator, hence all of $J$. Since $z_0N=0$,
$z_0x_G=z_0\rho(x_G)$. Apply $\ell$ to obtain the second
inclusion from the first.

Write $\lambda(G)=\rho(x_G)$. Define
\[
 \beta(G,H)=\ell(x_G)\lambda(H)\pmod {u^{-2b}V}
                         \quad\hbox{in }F/u^{-2b}V.
                                                               \tag{B45.3}
\]
This additive pairing is $A$-balanced. If $a\in A$ and
$a_0=\operatorname{pr}_0(a)$, the difference between
$\beta(aG,H)$ and $\beta(G,aH)$ equals
\[
 \ell\bigl((a_0-\rho(a_0))x_G\bigr)\lambda(H).
\]
The first factor lies in $u^{-b}V$ by (B45.2), since
$a_0-\rho(a_0)\in J$, and the second does also.
Their product therefore vanishes in the indicated quotient.
The quotient is an $A$-module through $\rho\operatorname{pr}_0$,
whose image is in $V$.

Choose an integer $r>2c+2b$. We construct a polynomial
$g_r\in F[X]$ with
\[
 g_r(0)=0,\quad g_r'(0)=(-u^{-1})^r,\quad
 g_r(V)\subset V,\quad g_r(O)\subset V+u^{-r}J.             \tag{B45.4}
\]
In characteristic $p$, choose a power $D$ of the residue
cardinality with $D\ge h$ and set
$\delta(X)=(X^D-X)/u$, $g_r=\delta^{\circ r}$.
The Frobenius identity and $N^h=0$ give
$\delta(a+n)=\delta(a)-n/u$ for $a\in V,n\in N$;
all assertions follow.

In mixed characteristic choose a power $D$ of the residue
cardinality so large that $D\ge h$ and
\[
 v(D)\ge (r-1)(h-2)+\max_{1\le j<h}v(j).
\]
Again use $g_r=\delta^{\circ r}$. The elementary identity
$\binom Dj=(D/j)\binom{D-1}{j-1}$ gives
$v(\binom Dj)\ge v(D)-v(j)$.
For $0\le i<r$, $a\in V$ and $n\in u^{-i}J$, expand
\[
 \delta(a+n)=\delta(a)+\frac{Da^{D-1}-1}{u}n
          +\sum_{2\le j<h}\frac{\binom Dj a^{D-j}}u n^j.
\]
The linear term belongs to $u^{-(i+1)}J$, and so does
each higher term: its valuation bound is
$v(\binom Dj)-ij-1\ge-i-1$.
Thus $\delta(V+u^{-i}J)\subset V+u^{-(i+1)}J$.
Induction proves (B45.4); its derivative follows from
$\delta(0)=0$ and $\delta'(0)=-u^{-1}$.

Extend polynomials from $Q_0$ to $Q(A)$ by zero on the
other factors, and put
\[
       F_r=u^c e z_0g_r,\qquad H_r=u^{r+c}eg_r.
\]
Both are in $\operatorname{Int}(A)$. For $H_r$ this is
(B45.1) and (B45.4). For $F_r$, multiplication by $z_0$
kills the nilpotent part of $g_r(a)$, leaving
$u^cz_0g_r(\rho(a))\in u^cO$.
The polynomial image of
\[
                   F_r\otimes H_r-H_r\otimes F_r
\]
is zero, because $F_r=u^{-r}z_0H_r$ in $Q(A)[X]$.
Its value under (B45.3) is $u^{2c-r}$: the second term
vanishes since $\rho(z_0)=0$, and the two signs cancel
in the first. This is nonzero modulo $u^{-2b}V$ by the
choice of $r$. The required tensor kernel is proved.

The structural inputs were opened in Stacks, Theorem 10.160.8
(Cohen structure, tag 032A), Lemma 10.160.2 (complete finite
algebras), and Lemma 10.162.8 (complete local rings are Nagata,
tag 032W). The lattice and obstruction arguments above are
given in full and do not invoke flatness of integer-valued
polynomial rings.
\end{proof}

\begin{saresult}{B.046}{Proved}{Complete one-dimensional fraction-field classification without finite normalization}
Let $R$ be an infinite one-dimensional Noetherian domain,
and put $K=\operatorname{Frac}(R)$.
The compatible coaddition and comultiplication tensors for
$\operatorname{Int}(R)$ are uniquely determined if and only
if, at each maximal ideal with finite residue field,
the completion $\widehat{R_{\mathfrak m}}$ is a domain.
Equivalently, at each such ideal the normalization of
$R_{\mathfrak m}$ is finite and local. In every positive
case all multivariable tensor maps are isomorphisms and
the compatible plethory is unique.
\end{saresult}
\begin{proof}
By B.033 the tensor question localizes, so fix a
finite-residue maximal ideal and write its local ring
as $R$. If its completion is nonreduced, B.045 applies:
the completion is one-dimensional Cohen--Macaulay, and
in mixed characteristic the nonzero element $p\in R$
remains a nonzerodivisor by flatness. B.044 then proves
failure of tensor injectivity over $R$.

If the completion is reduced, $R$ is analytically
unramified and its normalization $T$ is finite
(Stacks, Lemma 10.162.10(4), tag 032Y, whose full proof
was opened). B.025 proves that the tensors work exactly
when $T$ is local. This is equivalent to the completion
of $R$ being a domain. In one direction, when $T$ is
local it is a discrete valuation ring, and exact
completion embeds $\widehat R$ in the complete discrete
valuation ring $\widehat T$. Conversely, if
$\widehat R$ is a domain, flatness embeds
$T\otimes_R\widehat R$ into
$K\otimes_R\widehat R$, which is a field. But the
completion of the finite semilocal algebra $T$ is the
product of its nonzero completions at its maximal
ideals. A subring of a field cannot have the nontrivial
idempotents of a product with more than one factor.
Thus $T$ is local.

At maximal ideals of infinite residue field the ordinary
Vandermonde argument gives the polynomial ring, as
already proved. In the positive finite-residue case,
B.024 gives every multivariable tensor isomorphism.
Testing kernels and cokernels at every maximal ideal
proves the global assertion, and B.009 supplies the
unique plethory compatibility.
\end{proof}
